\section{数据模型与接口设计}

本章说明平台在 PostgreSQL + PostGIS 数据库下的数据模型、字段结构和接口规范。数据设计围绕“地图找点、规则判断、个性化推荐、出钓记录、社区分享、教程学习、活动服务”七类业务展开，要求结构清晰、字段可验证、空间查询可计算、隐私权限可控制。

\subsection{数据库设计原则}

数据库采用 PostgreSQL 作为主数据库，启用 PostGIS 扩展处理地理空间数据。PostgreSQL 负责保存用户、钓点、鱼种、出钓记录、帖子、教程、活动、订单、报告等结构化数据；PostGIS 负责钓点坐标、水域边界、禁钓区域、附近检索和空间叠加判断。

数据库设计遵循以下原则：

\begin{enumerate}
  \item \textbf{主数据稳定：}用户、钓点、鱼种、教程、装备等相对稳定数据采用独立主表维护，避免在记录和帖子中重复保存大量冗余信息。
  \item \textbf{业务数据可追溯：}出钓记录、社区帖子、评论、点赞、报告快照等业务数据保留创建时间、更新时间、操作者和状态字段，支撑统计分析与审计追踪。
  \item \textbf{空间数据可计算：}钓点、水域和禁钓区使用 PostGIS 几何字段保存，支持附近点位查询、距离排序、点面叠加和空间范围过滤。
  \item \textbf{隐私字段分级：}用户精确位置、私密记录、私信内容、图片地址等敏感字段与公开展示字段分开处理，接口返回时按权限和位置可见性进行脱敏。
  \item \textbf{扩展字段可控：}对于推荐结果、报告快照、第三方返回摘要等半结构化内容，采用 \texttt{jsonb} 保存，但核心关联关系仍通过关系表和外键表达。
\end{enumerate}

\subsection{核心数据表设计}

数据库表按业务域划分为用户与关系、垂钓主数据、记录与社区、推荐与报告、活动与服务、系统配置六类。

{\small
\begin{longtable}{|>{\raggedright\arraybackslash}p{3.9cm}|>{\raggedright\arraybackslash}p{3.0cm}|>{\raggedright\arraybackslash}p{6.7cm}|}
  \caption{数据库核心表设计}
  \label{tab:postgres-tables}\\
  \toprule
  \textbf{表名} & \textbf{用途} & \textbf{主要字段设计} \\
  \midrule
  \endfirsthead
  \multicolumn{3}{c}{\tablename~\thetable\quad（续）}\\
  \toprule
  \textbf{表名} & \textbf{用途} & \textbf{主要字段设计} \\
  \midrule
  \endhead
  \bottomrule
  \endlastfoot
  users & 用户基础信息 & id、nickname、avatar、phone\_hash、city、level、status、created\_at、updated\_at \\
  \hline
  user\_preferences & 用户个性化偏好 & user\_id、target\_fish、preferred\_methods、distance\_limit\_m、privacy\_default、risk\_tolerance、updated\_at \\
  \hline
  memberships & 会员状态与权益 & user\_id、member\_level、valid\_until、benefits、status、created\_at \\
  \hline
  fishing\_pois & 垂钓 POI 主表 & id、name、type、city、district、address、lng、lat、geom、fish、tags、score、facilities、is\_banned、compliance\_note、updated\_at \\
  \hline
  fishing\_areas & 水域与禁钓区域 & id、name、area\_type、city、geom、rule\_text、rule\_source、effective\_from、effective\_to、status \\
  \hline
  fish\_species & 鱼种知识库 & id、name、alias、habitat、active\_temperature、common\_methods、bait\_tips、season\_tips \\
  \hline
  weather\_snapshots & 天气快照 & id、city、adcode、weather、temperature\_c、wind\_level、humidity、pressure\_hpa、alerts、source、updated\_at \\
  \hline
  fishing\_records & 出钓记录 & id、user\_id、poi\_id、start\_time、end\_time、duration\_seconds、fish\_count、fish\_weight、method、bait、is\_blank\_trip、privacy\_level、location\_geom \\
  \hline
  catch\_entries & 单尾鱼获明细 & id、record\_id、species、weight、length\_cm、caught\_at、image\_url、note \\
  \hline
  posts & 社区帖子 & id、user\_id、post\_type、title、content、poi\_id、record\_id、visibility、location\_visibility、location\_text、display\_geom、created\_at \\
  \hline
  comments & 评论数据 & id、post\_id、user\_id、content、parent\_id、status、created\_at \\
  \hline
  reactions & 点赞、收藏等互动 & id、target\_type、target\_id、user\_id、reaction\_type、created\_at \\
  \hline
  recommendation\_logs & 推荐结果日志 & id、user\_id、request\_json、result\_json、agent\_version、created\_at \\
  \hline
  report\_snapshots & 个人报告快照 & id、user\_id、period、summary\_json、generated\_at、source \\
  \hline
  tutorials & 教程内容 & id、title、category、difficulty、format、duration、content、related\_poi\_ids、status \\
  \hline
  learning\_progress & 学习进度 & id、user\_id、tutorial\_id、status、practice\_notes、updated\_at \\
  \hline
  events & 活动信息 & id、title、city、poi\_id、start\_time、capacity、fee、safety\_notice、status \\
  \hline
  event\_registrations & 活动报名 & id、event\_id、user\_id、status、checked\_in\_at、created\_at \\
  \hline
  groups、messages & 活动群与消息 & group\_id、user\_id、content、message\_type、read\_status、created\_at \\
  \hline
  equipment、orders & 装备与订单 & equipment\_id、user\_id、quantity、order\_status、pay\_status、created\_at \\
\end{longtable}
}

\subsection{字段约束与索引设计}

空间字段使用 WGS84 坐标系，钓点位置采用 \texttt{geometry(Point, 4326)}，水域和禁钓区域采用 \texttt{geometry(Polygon, 4326)} 或 \texttt{geometry(MultiPolygon, 4326)}。\texttt{fishing\_pois.geom}、\texttt{fishing\_areas.geom} 建立 GiST 索引，用于附近检索和点面叠加判断。用户、记录、帖子、活动等高频查询字段建立 B-tree 索引，例如 \texttt{user\_id}、\texttt{poi\_id}、\texttt{created\_at}、\texttt{status}、\texttt{post\_type}。

关键约束包括：用户 ID、POI ID、记录 ID 和帖子 ID 采用全局唯一标识；记录开始时间不得晚于结束时间；鱼获数量和重量不得为负；帖子位置展示字段必须与 \texttt{location\_visibility} 保持一致；订单、活动报名、学习进度等状态字段使用枚举约束，防止非法状态进入业务流程。

\subsection{核心数据模型}

出钓记录是连接地图、推荐、报告和社区的核心数据。记录中既保存用户填写的鱼种、重量、钓法、饵料，也保存与 POI、天气和位置相关的结构化字段，用于个性化推荐和生涯统计。

\begin{lstlisting}[caption={出钓记录数据模型示例},label={lst:record-schema}]
{
  "id": "rec_001",
  "user_id": "u_001",
  "start_time": "2026-07-01T06:30:00",
  "end_time": "2026-07-01T10:45:00",
  "duration_seconds": 15300,
  "poi_id": "poi_wh_001",
  "location_name": "东湖磨山水域",
  "latitude": 30.5501,
  "longitude": 114.4085,
  "weather": "多云",
  "temperature": 28,
  "fish_count": 12,
  "fish_weight": 3.8,
  "fish_species": "鲫鱼,鲤鱼",
  "fishing_method": "台钓",
  "bait": "商品饵",
  "is_blank_trip": false,
  "privacy_level": "private"
}
\end{lstlisting}

\subsection{API 接口总览}

后端采用 FastAPI 提供统一 \texttt{/api} 前缀，接口以 RESTful 风格组织。接口覆盖地图 POI、天气、推荐、AI 解释、出钓记录、社区帖子、教程、用户报告、活动和服务订单等模块。

\begin{longtable}{|p{0.7cm}|p{3.5cm}|p{2.5cm}|p{3.9cm}|p{2.0cm}|}
  \caption{API 接口总览表}
  \label{tab:api-list}\\
  \toprule
  \textbf{序号} & \textbf{接口路径} & \textbf{方法} & \textbf{功能} & \textbf{所属模块} \\
  \midrule
  \endfirsthead
  \multicolumn{5}{c}{\tablename~\thetable\quad（续）}\\
  \toprule
  \textbf{序号} & \textbf{接口路径} & \textbf{方法} & \textbf{功能} & \textbf{所属模块} \\
  \midrule
  \endhead
  \bottomrule
  \endlastfoot
  1 & \path|/api/health| & {\scriptsize GET} & 服务健康检查 & 元数据 \\
  \hline
  2 & \path|/api/amap/config| & {\scriptsize GET} & 返回前端加载高德 JS API 所需配置 & 地图配置 \\
  \hline
  3 & \path|/api/pois| & {\scriptsize GET} & 查询钓点列表，支持城市、关键词、距离等条件 & 地图 POI \\
  \hline
  4 & \path|/api/pois/{poi_id}| & {\scriptsize GET} & 查询钓点详情 & 地图 POI \\
  \hline
  5 & \path|/api/weather/current| & {\scriptsize GET} & 获取实时天气快照 & 天气服务 \\
  \hline
  6 & \path|/api/recommendations| & {\scriptsize POST} & 根据位置、偏好、天气和候选 POI 返回推荐排序 & 推荐模型 \\
  \hline
  7 & \path|/api/ai/fishing-advice| & {\scriptsize POST} & 根据推荐结果生成个性化出钓建议 & AI Agent \\
  \hline
  8 & \path|/api/records| & {\scriptsize GET/POST} & 查询和新增出钓记录 & 出钓记录 \\
  \hline
  9 & \path|/api/records/{record_id}| & {\scriptsize\begin{tabular}[t]{@{}l@{}}GET\\PATCH\\DELETE\end{tabular}} & 查询、更新和删除单条出钓记录 & 出钓记录 \\
  \hline
  10 & \path|/api/posts| & {\scriptsize GET/POST} & 查询和发布社区帖子 & 社区 \\
  \hline
  11 & \path|/api/feed| & {\scriptsize GET} & 获取社区信息流 & 社区 \\
  \hline
  12 & \path|/api/posts/{post_id}| & {\scriptsize GET} & 查询帖子详情 & 社区 \\
  \hline
  13 & \path|/api/posts/{post_id}/comments| & {\scriptsize GET/POST} & 查询或新增评论 & 社区互动 \\
  \hline
  14 & \path|/api/posts/{post_id}/reactions| & {\scriptsize POST} & 点赞、收藏等互动 & 社区互动 \\
  \hline
  15 & \path|/api/tutorials| & {\scriptsize GET} & 查询教程列表 & 教程 \\
  \hline
  16 & \path|/api/tutorials/{tutorial_id}/progress| & {\scriptsize PUT} & 更新学习进度 & 教程 \\
  \hline
  17 & \path|/api/users/{user_id}| & {\scriptsize GET} & 查询用户基础信息 & 用户 \\
  \hline
  18 & \path|/api/users/{user_id}/profile-summary| & {\scriptsize GET} & 获取个人统计摘要 & 个人报告 \\
  \hline
  19 & \path|/api/users/{user_id}/reports| & {\scriptsize GET/POST} & 查询或生成报告快照 & 个人报告 \\
  \hline
  20 & \path|/api/fish-species| & {\scriptsize GET} & 查询鱼种知识库 & 知识库 \\
\end{longtable}

\subsection{接口规范}

接口响应统一采用“业务数据 + 元信息”的结构。列表接口返回 \texttt{data} 数组和 \texttt{meta} 元信息；详情接口返回单个 \texttt{data} 对象；创建接口成功时返回新建对象并使用 HTTP 201 状态码；删除接口成功时返回 HTTP 204。错误响应统一包含错误码、错误信息和请求追踪 ID，便于前端展示和后端排查。

\begin{lstlisting}[caption={列表接口返回结构示例},label={lst:list-response}]
{
  "data": [
    {
      "id": "poi_wh_001",
      "name": "东湖磨山水域",
      "score": 86
    }
  ],
  "meta": {
    "total": 1,
    "page": 1,
    "page_size": 20
  }
}
\end{lstlisting}

分页参数统一采用 \texttt{page}、\texttt{page\_size}；筛选参数按模块扩展，例如 POI 接口支持 \texttt{city}、\texttt{keyword}、\texttt{lat}、\texttt{lng}、\texttt{radius\_m}，帖子接口支持 \texttt{channel}、\texttt{post\_type}，记录接口支持 \texttt{user\_id}。接口错误处理遵循 HTTP 状态码语义：400 表示请求参数错误，401 表示未认证，403 表示无权限，404 表示资源不存在，409 表示状态冲突，500 表示服务端异常。生产环境中错误响应不返回堆栈和敏感配置。


